The following criterion gives two intervals of derivative ratios for
which every interior level has positive dimension.

\begin{lemma}\label{lem:two-bands}
Let \(0<\lambda<1/2\), and let \(f,g\) be increasing \(C^1\)
functions on \([0,1]\) with positive derivatives.
If all ratios \(g'(t)/f'(s)\), \(s,t\in[0,1]\), are greater than
one and belong to a compact subinterval of either
\begin{equation}\label{eq:two-bands}
\left(\frac1\lambda-1-\lambda,\frac1{1-\lambda}\right)
\quad\text{or}\quad
\left(\frac1\lambda-1,\frac1{1-\lambda-\lambda^2}\right),
\end{equation}
then every strict interior level of \(f+g\) on \([0,1]^2\)
meets \(C_\lambda^2\) in a set of positive Hausdorff dimension.
\end{lemma}
\begin{proof}
Subtract the values at zero. Choose strict bounds
\(1<L\leq g'(t)/f'(s)\leq U\) in one of the intervals
\eqref{eq:two-bands}. The four child image hulls, ordered as
\(00,10,01,11\), are
\begin{gather*}
[0,f(\lambda)+g(\lambda)],\quad
[f(1-\lambda),f(1)+g(\lambda)],\\
[g(1-\lambda),f(\lambda)+g(1)],\quad
[f(1-\lambda)+g(1-\lambda),f(1)+g(1)].
\end{gather*}
Consecutive hulls overlap strictly: the outer overlaps follow from
\(L>1/\lambda-2\), and the middle one from
\(U<1/(1-2\lambda)\). A level in no two child interiors must
belong to the respective closed regions
\begin{gather*}
[0,f(1-\lambda)],\quad
[f(\lambda)+g(\lambda),g(1-\lambda)],\\
[f(1)+g(\lambda),f(1-\lambda)+g(1-\lambda)],\quad
[f(\lambda)+g(1),f(1)+g(1)].
\end{gather*}
Empty regions are discarded. Restrict to the continuing child,
subtract its lower-corner value, and divide by \(\lambda\).
The derivative bounds are preserved. The inequalities
\begin{align*}
f(\lambda)+g(\lambda)&>f(1-\lambda^2),\\
g(1-\lambda)&<f(1)+g(\lambda^2),\\
f(1-\lambda)&<f(\lambda^2)+g(\lambda)
\end{align*}
follow respectively from
\(L>1/\lambda-1-\lambda\),
\(U<1/(1-\lambda-\lambda^2)\), and the first inequality on \(L\).
They show that a continuing \(10\) must be followed by \(10\),
and that \(00\) cannot be followed by \(11\), as long as the level
avoids strict overlap. Simultaneous reflection gives the corresponding
statements for \(01\) and \(11\).

The constant \(10\) address has level \(f(1)\). This is outside
its closed region: in the first band \(g(1-\lambda)<f(1)\),
whereas in the second band \(f(\lambda)+g(\lambda)>f(1)\).
Reflection excludes the constant \(01\) address. Thus an infinite
address avoiding overlap can only be \((00)^\infty\) or
\((11)^\infty\), at the two extreme levels. Every strict interior
level reaches a strict overlap in finitely many steps.

We need uniform returns to obtain positive dimension. Normalize
initially by \(\min f'=1\); then
\(1\leq f'\leq U/L\) and \(L\leq g'\leq U\).
The closed family of Lipschitz pairs with these individual increment
bounds and cross-variable increment-ratio bounds \([L,U]\) is
compact in the uniform topology. Translated cylinder restrictions,
divided by their geometric scale, preserve this whole family;
their minimum first derivative need not equal one.
All preceding comparisons use increment inequalities and hence
remain valid on the closed family.

Choose \(0<\delta<\lambda(1+L)/(1+\lambda)\).
Use \([\delta,f(1)+g(1)-\delta]\) as the central level domain.
The images of central domains under all-zero prefixes are
\[
[\delta\lambda^j,
 f(\lambda^j)+g(\lambda^j)-\delta\lambda^j],\qquad j\geq0.
\]
Successive intervals overlap strictly by the lower increment bounds
and the choice of \(\delta\). Their interiors reach every positive
level below the upper endpoint of the original central domain;
reflection reaches the remaining interior levels. At a strict overlap,
both children can therefore be brought strictly into their central
domains by finite prefixes. These are two incomparable returning words.
The space of function pairs and their central levels is compact,
and each such pair of strict returns remains valid on an open
neighbourhood. A finite subcover gives a maximum return length
\(M\) and a positive return margin. Iteration produces a separated
binary tree of product cylinders. The cylinder separation
\(1-2\lambda>0\) and the mass distribution principle give dimension
at least \(\log2/(M|\log\lambda|)>0\). An arbitrary interior level
first enters a central domain. This initial finite restriction
does not change positivity of dimension.
\end{proof}

We apply this criterion to actual rectangles and to endpoint models.
Let \([u_-,u_+]\), \([v_-,v_+]\) be positive cylinder hulls of
depths \(p,q\). An actual rectangle satisfies
\begin{equation}\label{eq:uniform-actual}
u_-^\alpha+v_-^\alpha<1<u_+^\alpha+v_+^\alpha,
\end{equation}
and its derivative-ratio interval is
\[
\left[\lambda^{q-p}(v_-/u_+)^{\alpha-1},
      \lambda^{q-p}(v_+/u_-)^{\alpha-1}\right].
\]
An endpoint model, with integers \(j\) and \(k\geq0\), satisfies
\begin{equation}\label{eq:uniform-endpoint}
u_-^\alpha<\alpha\lambda^{j-k\alpha}v_+,
\qquad u_+^\alpha>\alpha\lambda^{j-k\alpha}v_-.
\end{equation}
After reflection of the second coordinate, the power form
\(u^\alpha-\alpha\lambda^{j-k\alpha}v\) has derivative-ratio interval
\[
\left[\lambda^{j-k\alpha+q-p}u_+^{1-\alpha},
      \lambda^{j-k\alpha+q-p}u_-^{1-\alpha}\right],
\]
whereas the zero-equivalent logarithmic form
\(\alpha\log u-\log v-\log \alpha-(j-k\alpha)\log\lambda\) has
first-to-second derivative-ratio interval
\begin{equation}\label{eq:uniform-log-ratio}
\left[\alpha\lambda^{p-q}v_-/u_+,
      \alpha\lambda^{p-q}v_+/u_-\right].
\end{equation}
In each case we may take the reciprocal interval by exchanging
coordinates. A source is accepted when its whole ratio interval
lies strictly above one and inside one band of \eqref{eq:two-bands}.

\begin{theorem}\label{thm:last-real}
For every real \(\alpha>1\),
\[
\lambda_{\max}(\alpha)<\frac{\sqrt3-1}{2}.
\]
\end{theorem}
\begin{proof}
Put $c:=(\sqrt3-1)/2$. We prove positive Hausdorff dimension for
every $c\le\lambda<1/2$, and obtain the strict inequality by preserving
a finite collection of strict returns when $\lambda$ varies near $c$.
The irrational exponents admit direct endpoint sources. For rational
exponents we first construct infinitely many intersections, then use
the arithmetic of isolated points to obtain a perfect subset and strict
returns. The latter argument uses Faltings's theorem and a theorem of
Adamczewski and Bugeaud on digit expansions.

\paragraph{Transfer from endpoint models.}
A rectangle satisfying \eqref{eq:uniform-endpoint}, with its entire
derivative-ratio interval strictly in one of \eqref{eq:two-bands} and
strictly above one, gives positive dimension for the actual curve.
To see this for any real $\alpha>1$, choose positive integers $t_i$ with
$\displaystyle\lim_{i\to\infty}t_i=\infty$ and integers $r_i$ with
$\displaystyle\lim_{i\to\infty}(r_i-\alpha t_i)=0$.
For rational $\alpha$ use multiples of its denominator; for irrational
$\alpha$ use recurrence modulo one. Set
\[
x:=\lambda^{k+t_i}u,\qquad y:=1-\lambda^{j+r_i}v.
\]
For large $i$ these are Cantor coordinates, and Taylor's formula gives
\[
\lim_{i\to\infty}
\lambda^{-\alpha(k+t_i)}(x^\alpha+y^\alpha-1)
=u^\alpha-\alpha\lambda^{j-k\alpha}v
\quad\text{in }C^1
\]
on the fixed positive rectangle. For a logarithmic source use instead
the exactly zero-equivalent expression
\[
\alpha\log u-
\log\left(\frac{1-(1-\lambda^{j+r_i}v)^\alpha}
{\lambda^{\alpha(k+t_i)}}\right).
\]
Its $C^1$ limit is the logarithmic endpoint expression used in
\eqref{eq:uniform-log-ratio}. Reflection of the second coordinate makes
both summands increasing. For a sufficiently large fixed $i$, all
strict corner and derivative inequalities persist, so
Lemma~\ref{lem:two-bands} gives positive dimension. The affine map back
to $(x,y)$ preserves dimension. These finitely many strict inequalities
also persist when $\lambda$ varies near its original value.

\par\medskip\noindent\emph{Irrational exponents.}
Let $\alpha>1$ be irrational and $c\le\lambda<1/2$. We construct a strict logarithmic endpoint source. First, for every $\lambda\ge1/3$, the outer Cantor cylinder satisfies
\[
\frac{(1-\lambda)+\lambda C_\lambda}
     {(1-\lambda)+\lambda C_\lambda}
=\left[1-\lambda,\frac1{1-\lambda}\right].
\]
To prove this, fix $1-\lambda\le r\le1$. The four child hulls for $s+rt$, in increasing order, are
\[
[0,\lambda(1+r)],\quad
[r(1-\lambda),r+\lambda],\quad
[1-\lambda,1+r\lambda],\quad
[(1+r)(1-\lambda),1+r].
\]
The two outer overlaps follow from $r(1-2\lambda)\le\lambda$, and the middle overlap from $r\ge1-2\lambda$. These inequalities hold because $r\le1$, $r\ge1-\lambda$ and $\lambda\ge1/3$. The same covering holds in every equal-depth descendant. Nested compactness therefore gives $C_\lambda+rC_\lambda=[0,1+r]$. Since $C_\lambda=1-C_\lambda$, it follows that $C_\lambda-rC_\lambda=[-r,1]$. The difference of the outer cylinder and $r$ times itself is consequently
\[
[1-\lambda-r,\,1-r(1-\lambda)],
\]
which contains zero. For $1\le r\le1/(1-\lambda)$, apply the preceding argument to $1/r$ and exchange the two factors. The hull endpoints exclude every other quotient, proving the displayed identity.

Consider the trimmed first slope band
\[
\left(
\frac{51}{50}\max\{1,1/\lambda-1-\lambda\},
\frac{50}{51(1-\lambda)}
\right).
\]
It is nonempty: the function
$\max\{1,1/\lambda-1-\lambda\}(1-\lambda)$ decreases with $\lambda$ and equals $\sqrt3/2$ at $c$, while $2601\sqrt3<5000$.
The quotient identity shows that every number in
\[
\left(
\frac{51}{50}\max\{1,1/\lambda-1-\lambda\}(1-\lambda),
\frac{50}{51(1-\lambda)^2}
\right)
\]
can be written as $ru/v$, where $r$ lies in the trimmed band and $u,v$ belong to the outer Cantor cylinder. This interval has multiplicative width greater than $1/\lambda$, because
\[
\lambda>
\left(\frac{51}{50}\right)^2
\max\{1,1/\lambda-1-\lambda\}(1-\lambda)^3.
\]
When the maximum is $1/\lambda-1-\lambda$, the right side decreases with $\lambda$, and at $c$ the difference between the two sides is
\[
\frac{18409-10606\sqrt3}{10000}>0,
\qquad 18409^2-3\cdot10606^2=1429573.
\]
When the maximum is one, its right side is less than
$(51/50)^2\,8/27<1/3<\lambda$.
The lower endpoint of the preceding interval is also less than one, since it is at most $51\sqrt3/100<1$.

Choose the first integer $n\ge0$ for which $\alpha\lambda^n$ is strictly below that interval's upper endpoint. If $n>0$, minimality puts it at least at $\lambda$ times the upper endpoint, hence strictly above the lower endpoint by the multiplicative-width inequality. If $n=0$, the lower comparison follows from $\alpha>1$. We may therefore choose outer Cantor points $u,v$ such that
\[
\frac{51}{50}\max\{1,1/\lambda-1-\lambda\}
<\alpha\lambda^n\frac vu
<\frac{50}{51(1-\lambda)}.
\]
Take cylinders containing $u$ and $v$ at depths $n+8$ and $8$, and denote their hulls by $[u_-,u_+]$ and $[v_-,v_+]$. They lie in the outer cylinder. Each relative endpoint error is at most
\[
\frac{\lambda^8}{1-\lambda}\le\frac1{128}.
\]
Thus every ratio in
\[
\left[\alpha\lambda^n\frac{v_-}{u_+},
      \alpha\lambda^n\frac{v_+}{u_-}\right]
\]
differs from $\alpha\lambda^n v/u$ by a multiplicative factor of at most $129/127<51/50$. The entire interval consequently lies strictly above one and strictly within the first band of Lemma~\ref{lem:two-bands}. After reflecting the $v$ coordinate and placing it first, this is precisely the derivative-ratio interval of the normalised function $\alpha\log u-\log v$ on the two cylinders.

It remains to choose an admissible endpoint coefficient. The open interval
\[
\left(
\alpha\log u_- -\log\alpha-\log v_+,
\alpha\log u_+ -\log\alpha-\log v_-
\right)
\]
has length
\[
\alpha\log\frac{u_+}{u_-}+\log\frac{v_+}{v_-}
\ge v_+-v_-=\lambda^8>0.
\]
The inequality uses $1/t\ge1$ on the positive cylinder in $[0,1]$. Since $\alpha$ is irrational, the set $\{j-k\alpha:j\in\mathbb Z,\ k\ge0\}$ is dense in $\mathbb R$. Indeed the pigeonhole principle supplies positive iterates of the rotation by $-\alpha$ arbitrarily close to zero modulo one; repeated multiples of each such iterate form arbitrarily fine nets of the circle, and integer translates give density on the line. Hence the displayed open interval contains $(j-k\alpha)\log\lambda$ for some integer $j$ and integer $k\ge0$.
This gives the strict corner inequalities
\[
u_-^\alpha<\alpha\lambda^{j-k\alpha}v_+,
\qquad u_+^\alpha>\alpha\lambda^{j-k\alpha}v_-.
\]
The logarithmic endpoint-to-actual transfer established above now supplies an actual Fermat source to which Lemma~\ref{lem:two-bands} applies. No assumption that $u$ or $v$ is an interior point of its cylinder was needed. At $\lambda=c$, the chosen words and integers give finitely many strict inequalities, which remain valid for $\lambda$ in a neighbourhood of $c$. This proves the required strict bound for every irrational $\alpha>1$.

% Exceptional exponents in the proof of Theorem 4.2.
\paragraph{The exponents $2$, $3$, and $3/2$.}
These three exponents admit explicit sources for the endpoint model
\eqref{eq:uniform-endpoint}, with $j,k$ as defined there.
A binary word specifies the cylinder obtained by its successive left or
right branches. We verify the strict corner inequalities and place each
whole derivative-ratio interval strictly above one inside a band of
\eqref{eq:two-bands}. All estimates below hold for $c\le\lambda<1/2$;
we use $c^2+c=1/2$ and $4/11<c<11/30$.

For $\alpha=3$, take $(j,k)=(2,0)$, so the endpoint equation is
$u^3=3\lambda^2v$.
For $c\le\lambda\le3/8$, use the words $1010$ and $11$;
for $3/8\le\lambda<1/2$, use $101$ and $1$.
In both cases
\[
 u_-=1-\lambda+\lambda^2-\lambda^3,
\]
while $u_+$ is respectively
$1-\lambda+\lambda^2-\lambda^3+\lambda^4$ and $1-\lambda+\lambda^2$.
Both depth differences are two, so the ratio interval is
$[u_+^{-2},u_-^{-2}]$. We place it in the second band.
For its upper bound,
\[
 (1-\lambda+\lambda^2-\lambda^3)^2-(1-\lambda-\lambda^2)
 >\lambda[-1+4\lambda(1-\lambda)+2\lambda^3]>0.
\]
The first inequality discards the positive term $(1-\lambda)^2\lambda^4$.
The bracket increases up to $1/2$ and at $c$ equals $11c-4>0$.
For the first source, $u_+$ decreases on $[0,1/2]$, since its derivative is
$(2\lambda-1)+\lambda^2(4\lambda-3)<0$.
Its value at $4/11$ is $10805/14641<3/4$, whence
\[
 (1-\lambda)u_+^2<\frac7{11}\left(\frac34\right)^2
 =\frac{63}{176}<\frac4{11}<\lambda.
\]
For the second source, both factors decrease, and
\[
 (1-\lambda)u_+^2\le\frac58\left(\frac{49}{64}\right)^2
 =\frac{12005}{32768}<\frac38\le\lambda.
\]
These are exactly the lower-band comparisons, and also show that the ratios
exceed one.
The lower corner holds for both sources: $u_-$ decreases and its value
at $4/11$ is $959/1331<8/11$, so
\[
 u_-^3<\frac{512}{1331}<\frac{528}{1331}
 =3\left(\frac4{11}\right)^2<3\lambda^2=3\lambda^2v_+.
\]
For the first upper corner, $u_+$ decreases while
$3\lambda^2(1-\lambda^2)$ increases on $[c,3/8]$.
At $3/8$, $u_+=3001/4096>23/32$, and
\[
 \left(\frac{23}{32}\right)^3=\frac{12167}{32768}
 >\frac{11880}{32768}
 =3\left(\frac38\right)^2\left(1-\left(\frac38\right)^2\right).
\]
For the second upper corner,
\[
 u_+^3\ge\left(\frac34\right)^3=\frac{27}{64}
 >\frac38\ge3\lambda^2(1-\lambda)=3\lambda^2v_-.
\]
The last bound uses the increase of $\lambda^2(1-\lambda)$ up to $1/2$.

For $\alpha=2$, take $(j,k)=(1,0)$, giving $u^2=2\lambda v$.
Use the words $1001,100$ on $[c,147/400]$;
$1010,101$ on $[147/400,3/8]$; and $101,10$ on $[3/8,1/2)$.
In every case the first depth exceeds the second by one, so the ratio
interval is $[1/u_+,1/u_-]$. Since $u_->1-\lambda$ and $u_+<1$,
its upper endpoint is strictly below $1/(1-\lambda)$ and the interval
lies above one. The remaining first-band comparison is
\[
 u_+(1-\lambda-\lambda^2)<\lambda.
\]
For the first source, $u_+=1-\lambda+\lambda^3$ decreases, as does its
product with $1-\lambda-\lambda^2$. At $c$ this product equals
$(1+c)/4<c$, because $c>1/3$.
For the middle source,
$u_+=1-\lambda+\lambda^2-\lambda^3+\lambda^4$ decreases as above.
At $c$ it equals $11/4-(11/2)c<737/1000<59/80$, using $c>183/500$.
Hence, for $\lambda\ge147/400$,
\[
 u_+(1-\lambda-\lambda^2)
 <\frac{59}{80}\frac{79591}{160000}<\frac{147}{400}\le\lambda,
\]
where $59\cdot79591=4695869<4704000=147\cdot32000$.
For the final source,
\[
 u_+(1-\lambda-\lambda^2)
 \le\frac{49}{64}\frac{31}{64}
 =\frac{1519}{4096}<\frac38\le\lambda.
\]

For the first source's lower corner,
$u_-=1-\lambda+\lambda^3-\lambda^4$ decreases, and at $c$ it equals
$(5/2)c-1/4<2/3$, since $c<11/30$.
Also $v_+=1-\lambda+\lambda^3\ge347/512>2/3$ on this interval.
Thus $u_-^2<4/9<2\lambda v_+$.
For its upper corner, $u_+$ decreases and $2\lambda(1-\lambda)$ increases,
so it suffices to compare at $147/400$.
There $u_+>341/500$, because
$147^3=3176523>3168000=99\cdot32000$, and
\[
 \left(\frac{341}{500}\right)^2
 -2\frac{147}{400}\frac{253}{400}
 =\frac{473}{2000000}>0.
\]
For the other two lower corners, the common endpoints satisfy
\[
 u_-=1-\lambda+\lambda^2-\lambda^3<\frac8{11},\qquad
 v_+=1-\lambda+\lambda^2\ge\frac34,
\]
so $u_-^2<64/121<6/11<2\lambda v_+$.
For the middle upper corner, $u_+$ decreases while
$\lambda(1-\lambda+\lambda^2-\lambda^3)$ increases, since its derivative is
$(1-2\lambda)+\lambda^2(3-4\lambda)>0$.
At $3/8$,
\[
 u_+^2=\left(\frac{3001}{4096}\right)^2
 >\frac{1095}{2048}=2\frac38\frac{365}{512}=2\lambda v_-;
\]
the integer comparison is $3001^2=9006001>8970240=1095\cdot8192$.
For the final upper corner,
\[
 u_+^2\ge\left(\frac34\right)^2=\frac9{16}
 >\frac12\ge2\lambda(1-\lambda)=2\lambda v_-.
\]

For $\alpha=3/2$, one source works throughout the interval.
Use the words $101$ and $10$, whose hulls are
\[
 [1-\lambda+\lambda^2-\lambda^3,1-\lambda+\lambda^2],
 \qquad[1-\lambda,1-\lambda+\lambda^2].
\]
Take $(j,k)=(2,1)$, giving
$u^{3/2}=(3/2)\lambda^{1/2}v$.
The ratio interval is
\[
 \left[(\lambda(1-\lambda+\lambda^2))^{-1/2},
       (\lambda(1-\lambda+\lambda^2-\lambda^3))^{-1/2}\right].
\]
Its lower endpoint lies above $1/\lambda-1$ precisely when
\[
 \lambda>(1-\lambda)^2(1-\lambda+\lambda^2).
\]
The right side decreases, and at $c$ the difference is
$(58c-21)/4>1/44>0$, using $c>4/11$.
The upper endpoint lies below $1/(1-\lambda-\lambda^2)$ precisely when
\[
 (1-\lambda-\lambda^2)^2
 <\lambda(1-\lambda+\lambda^2-\lambda^3).
\]
The left side decreases, whereas the right side increases by the derivative
calculation above. At $c$ they equal $1/4$ and $(22c-7)/4$, respectively;
the latter is greater because $c>4/11$.
Thus the entire ratio interval is strictly inside the second band and
above one.
For the lower corner, the decreasing endpoint satisfies
\[
 u_-<1-\frac13+\left(\frac13\right)^2-\left(\frac13\right)^3
 =\frac{20}{27}<\frac34<\frac94\lambda.
\]
Since $u_-<u_+=v_+$, we obtain
$u_-^3<u_-u_+^2<(9/4)\lambda v_+^2$.
For the upper corner, elementary differentiation gives
$\lambda(1-\lambda)^2\le4/27$, so
\[
 u_+^3\ge\left(\frac34\right)^3=\frac{27}{64}
 >\frac13\ge\frac94\lambda(1-\lambda)^2.
\]
Taking square roots proves both required corner inequalities.

By the preceding endpoint-to-actual transfer and
Lemma~\ref{lem:two-bands}, these sources give positive-dimensional actual
intersections for every $c\le\lambda<1/2$.
For each exponent, the source used at $\lambda=c$ has strict margins,
so one fixed sufficiently deep actual source persists on a neighbourhood
of $c$. Together these facts give the required strict last-transition
bound for $\alpha=2,3,3/2$.

% Endpoint construction for the proof of Theorem 4.2.
% Assumes c=(sqrt(3)-1)/2 and c<=lambda<1/2.
\paragraph{Nonempty endpoint fibres.}
We first construct infinitely many nontrivial points of the intersection for
rational $\alpha>1$. The construction will also give the strict margins needed
below. Put
\[
 A:=\log\bigl((1-\lambda)+\lambda C_\lambda\bigr),
 \qquad b:=-\log(1-\lambda),\qquad h:=-\log\lambda.
\]
Thus $\operatorname{conv}A=[-b,0]$ and $0<b<h$.
We will show that, for every real $\alpha>1$, there is an integer $j$ such that
\begin{equation}\label{eq:coarse-strict-link}
 \alpha A\cap
 (\log\alpha-jh-b,\log\alpha-jh)\ne\varnothing.
\end{equation}
This is a statement about a point lying strictly inside a translated hull;
we will use the Gap Lemma to obtain actual intersections of Cantor sets.

First, $A$ has thickness greater than one. Consider a cylinder hull $[a,a+d]$
inside $[1-\lambda,1]$, where $d=\lambda^n$ and $n\ge1$.
Its central gap is $(a+\lambda d,a+(1-\lambda)d)$.
Since the derivative of the logarithm decreases, the logarithmic left bridge
is at least $\lambda/(1-2\lambda)>1$ times the logarithmic gap.
The right bridge is longer than the gap precisely when
\[
 (a+d)(a+\lambda d)>(a+(1-\lambda)d)^2,
\]
or, equivalently,
\[
 (3\lambda-1)+(3\lambda-1-\lambda^2)\frac da>0.
\]
Here $0<d/a\le\lambda/(1-\lambda)$. If the second coefficient is negative,
the left side is at least
\[
 \frac{3\lambda-1-\lambda^3}{1-\lambda}>0;
\]
the numerator increases on $[c,1/2]$ and its value at $c$ is $(3c-1)/2>0$.
If the coefficient is nonnegative, positivity is immediate.
The bridge-to-gap ratios extend continuously to $d/a=0$, with limiting
value $\lambda/(1-2\lambda)>1$, so they have a strict common lower bound
above one for each fixed $\lambda$.
Moreover, every logarithmic gap one generation deeper is at most
$\lambda/(1-\lambda)<1$ times every gap of the preceding generation:
this follows by integrating $1/x$ over the corresponding gaps in
$[1-\lambda,1]$. Hence later gaps cannot truncate these bridges when the
gaps are removed in decreasing order. This proves the thickness assertion.

The Gap Lemma~\cite{Newhouse1979} now gives the exact identity
\begin{equation}\label{eq:coarse-thickening}
 \alpha A-A=\alpha A+[0,b],\qquad \alpha\ge1.
\end{equation}
Indeed, if $s\in\alpha A+[0,b]$, the hull of $s+A$ meets $\alpha A$.
A contact at a hull endpoint already gives an intersection.
Otherwise $s+A$ is not contained in a gap of $\alpha A$, because its hull
meets $\alpha A$, and $\alpha A$ is not contained in a gap of $s+A$,
because its diameter is at least $b$, larger than every bounded gap of $A$;
the hull intersection excludes containment in an unbounded gap.
The thickness product exceeds one, so the two Cantor sets intersect.
The reverse inclusion is immediate from $-A\subset[0,b]$.
The largest gap of $A$ has length
\[
 \log\frac{1-\lambda^2}{1-\lambda+\lambda^2}.
\]
Consequently, if this length multiplied by $\alpha$ is strictly less than
$b$, then $\alpha A+[0,b]=[-\alpha b,b]$, with strict overlap across each
gap. If also $(\alpha+1)b>h$, every residue modulo $h$ has a representative
$s$ for which $\alpha A\cap(s-b,s)\ne\varnothing$.
More generally, the same conclusion holds whenever every complementary arc
of $\alpha A$ modulo $h$ has length strictly less than $b$.

For $1<\alpha\le23/19$, we have
$0<\log\alpha<b$, since $23/19<3/2<1/(1-c)$.
The point $0\in\alpha A$ therefore proves
\eqref{eq:coarse-strict-link} with $j=0$.
For $23/19\le\alpha\le3$, the ratio
$(1-\lambda^2)/(1-\lambda+\lambda^2)$ decreases with $\lambda$:
its derivative has the sign of $1-4\lambda+\lambda^2<0$.
Using $\lambda\ge c>4/11$ gives
\[
 \left(\frac{1-\lambda^2}{1-\lambda+\lambda^2}\right)^3(1-\lambda)
 <\left(\frac{35}{31}\right)^3\frac7{11}
 =\frac{300125}{327701}<1.
\]
Thus all gaps of $\alpha A$ are shorter than $b$.
Also $(42/19)b>h$. It suffices to verify this at $\lambda=c$;
using $1-c=\sqrt3c$, it is equivalent to $3^{21}c^{23}<1$.
Since $c<11/30$ and
\[
 \left(\frac{11}{10}\right)^{10}<\frac{13}{5},\qquad
 \left(\frac{11}{10}\right)^{23}
 <\left(\frac{13}{5}\right)^2\left(\frac{11}{10}\right)^3
 =\frac{224939}{25000}<9,
\]
where the first comparison is $25937424601<26000000000$, the desired
inequality follows. This proves \eqref{eq:coarse-strict-link} up to $\alpha=3$.

We next prove that every complementary circle arc is shorter than $b$ for
$3\le\alpha\le6$. Reflecting $A$, write its root gap in $[0,b]$ as $(L,U)$,
where
\[
 L:=-\log(1-\lambda^2),\qquad U:=-\log(1-\lambda+\lambda^2).
\]
Every other logarithmic gap has length at most
$\lambda^2(1-2\lambda)/(1-\lambda)<2/27$;
here $\lambda^2(1-2\lambda)\le1/27$ and $1-\lambda>1/2$.
Their dilations by $\alpha\le6$ therefore have length less than $4/9<b$.
The elementary bounds
\begin{equation}\label{eq:coarse-basic-bounds}
 b>\frac{91}{200},\qquad h<\frac{503}{500}
\end{equation}
will be justified explicitly below. Furthermore,
\[
 b-U=\log\left(1+\frac{\lambda^2}{1-\lambda}\right)>
 \frac{16}{93},\qquad
 U-L<\log\frac{35}{31}<\frac4{31},\qquad
 h-b<\log\frac74<\frac35.
\]
The first bound follows from $\lambda>4/11$ and
$\log(1+x)>x/(1+x)$; the last follows from
$1+3/5+(3/5)^2/2=89/50>7/4$. Hence
\[
 b(b-U)>\frac{182}{2325}>\frac{12}{155}>(U-L)(h-b),
 \qquad \frac{182}{2325}-\frac{12}{155}=\frac{62}{72075}>0.
\]
Also $\alpha b>h$ and $\alpha(U-L)<24/31<h+b$, using
$h>\log2>1/2$.
If an empty circle arc had length at least $b$, choose an empty open subarc
$I=(a,a+b)\subset(0,h)$, allowing endpoints $0,h$.
It lies in the root gap $(\alpha L,\alpha U)$, so $\alpha(U-L)\ge b$.
Its translate $I+h$ cannot partly meet $[0,\alpha b]$, since the hull endpoint
$\alpha b$ belongs to $-\alpha A$; nor can it lie wholly in that hull,
since then the same root gap would have length at least $h+b$.
Thus $a\ge\alpha b-h$, and $a+b\le\alpha U$ gives
$\alpha(b-U)\le h-b$. The last two inequalities contradict
$b(b-U)>(U-L)(h-b)$.

For $6\le\alpha\le20$, only three logarithmic gaps can be relevant.
Their lengths, in decreasing generations, are
\[
 \log\frac{1-\lambda^2}{1-\lambda+\lambda^2},\qquad
 \log\frac{(1-\lambda)(1+\lambda^2)}{1-\lambda+\lambda^3},\qquad
 \log\frac{1-\lambda^3}{1-\lambda^2+\lambda^3}.
\]
They are respectively less than
\begin{equation}\label{eq:coarse-three-gaps}
 \frac{33}{125}b,\qquad\frac{113}{1000}b,\qquad\frac{17}{200}b;
 \qquad \frac bh>\frac{453}{1000},
\end{equation}
and every deeper gap is shorter than $b/20$.
Here are explicit bounds for these claims and for
\eqref{eq:coarse-basic-bounds}. We have $c>183/500$.
The derivatives of the three logarithmic lengths have, after multiplication
by positive denominators, the numerators
\[
 1-4\lambda+\lambda^2,\quad
 \lambda(2-7\lambda+4\lambda^2-\lambda^3),\quad
 \lambda(2-6\lambda+\lambda^3).
\]
Each is negative for $183/500\le\lambda\le1/2$, since the relevant polynomial
is negative at $183/500$ and decreases thereafter. The logarithmic lengths
therefore decrease while $b$ increases.
The positive logarithmic series gives
\[
 b>2\left(\frac{183}{817}+\frac{(183/817)^3}{3}
                   +\frac{(183/817)^5}{5}\right)>\frac{45569}{100000}.
\]
For $0<z<1$,
\[
 \log\frac{1+z}{1-z}
 \le2\left(z+\frac{z^3}{3}+\frac{z^5}{5(1-z^2)}\right).
\]
Substituting $z=(R-1)/(R+1)$ gives the following rational bounds:
\[
\begin{array}{c|c}
R&2\bigl(z+z^3/3+z^5/[5(1-z^2)]\bigr)\ \text{is less than}\\ \hline
216511/191989&12021/100000\\
89866013/85378487&5123/100000\\
118871513/114383987&3849/100000
\end{array}
\]
The left column gives the three gap ratios at $\lambda=183/500$.
The entries in the right column are respectively smaller than
$33/125,113/1000,17/200$ times $45569/100000$.
Also
\[
 \sum_{k=0}^{8}\frac{(10053/10000)^k}{k!}-\frac{500}{183}
 =\frac{39624613635943016066263577308071007}
 {81984000000000000000000000000000000000}>0.
\]
Thus $h<10053/10000<503/500$, and
$(45569/100000)/(10053/10000)>453/1000$.
For the largest depth-three gap, $\log(1+x)\le x$ and
$b\ge\lambda+\lambda^2/2$ give the ratio bound
\[
 \frac{\text{gap length}}b
 \le\frac{(1-2\lambda)\lambda^2}{(1-\lambda)(1+\lambda/2)}<\frac1{20}.
\]
The rational function decreases on $[183/500,1/2]$, because its derivative
has the sign of $2-(13/2)\lambda+2\lambda^2+\lambda^3<0$.
At $183/500$ its difference from $1/20$ is
$400223/187505500>0$. The ordering by generation treats all later gaps.

Suppose now that an open circle arc of length $b$ misses $\alpha A$.
At least $\lfloor(\alpha-1)b/h\rfloor$ of its translates lie wholly in
$[-\alpha b,0]$. Each must lie in a gap of $\alpha A$.
A gap of length $\alpha g$ accommodates at most
$\lfloor(\alpha g-b)/h\rfloor+1$ such translates, and none if $\alpha g<b$.
Use the normalized hull length $t:=b/h\in(453/1000,1)$.
If $k$ of the three upper bounds in \eqref{eq:coarse-three-gaps} can be active,
the total capacity is at most
$\alpha t$ times their sum, minus $kt$, plus $k$.
The required number is greater than $(\alpha-1)t-1$.
Thus capacity is insufficient whenever
\[
 [\alpha(1-\text{sum of the active upper ratios})+k-1]t>k+1.
\]
For $6\le\alpha<1000/113$, only the first gap can be active, and
$6(1-33/125)(453/1000)-2=7/15625>0$.
For $1000/113\le\alpha<200/17$, the sum is $377/1000$.
The sufficient inequality holds if $\alpha\ge91/10$ or $t\ge47/100$;
the respective endpoint margins are
$211929/10000000$ and $173/2825$.
In the remaining rectangle, each capacity is at most one, whereas
$\lfloor(\alpha-1)t\rfloor\ge3$.
For $200/17\le\alpha\le20$, the sum is $231/500$.
The sufficient inequality holds if $\alpha\ge13$ or $t\ge49/100$;
the endpoint margins are $37141/500000$ and $173/2125$.
In the remaining rectangle the two smaller capacities are at most one
and the root capacity is at most two. If the latter is at most one,
the total is at most three, but the required number is at least four.
If it is two, then $((33/125)\alpha-1)t\ge1$.
For $\alpha\le25/2$ this gives
\[
 (\alpha-1)t\ge\frac{\alpha-1}{(33/125)\alpha-1}\ge5;
\]
for $\alpha\ge25/2$ the bound $t>453/1000$ gives the same conclusion.
The required number is therefore at least five, against capacity at most four.
All cases are impossible, proving the strict circle-gap bound.

It remains to treat $\alpha\ge20$.
For $1\le\beta<1/\lambda$, the largest complementary arc of
$\beta C_\lambda$ modulo one is at most
\[
 \begin{cases}
 \beta(1-2\lambda),&1\le\beta\le1/(1-\lambda),\\
 2-\beta,&1/(1-\lambda)\le\beta\le1+2\lambda,\\
 1-2\lambda,&1+2\lambda\le\beta<1/\lambda.
 \end{cases}
\]
Indeed, all nonroot gaps have length at most $1-2\lambda$.
A longer empty arc has a representative in $(0,1)$ lying in the root gap.
Its translate by one cannot partly meet $[0,\beta]$, because $\beta$ belongs
to the Cantor set, and cannot lie wholly in that hull, because the root gap
has length $\beta(1-2\lambda)<1$. It therefore lies outside the hull.
The first arc is consequently contained in
\[
 (\max(\beta\lambda,\beta-1),\min(\beta(1-\lambda),1)).
\]
Taking the maximum of this interval length and $1-2\lambda$ gives the
three displayed pieces. This bound decreases with $\lambda$ for fixed $\beta$.

Choose the largest integer $n\ge1$ such that
\[
 \beta:=\frac{\alpha}{h}
       \log\left(1+\frac{\lambda^n}{1-\lambda}\right)\ge1.
\]
Such an $n$ exists because the value at $n=1$ is $\alpha b/h>1$.
Concavity of the logarithm and maximality give $\beta<1/\lambda<11/4$.
The set $\alpha A$ contains the image
$\alpha\log((1-\lambda)+\lambda^n C_\lambda)$.
Compared with its secant parametrisation
$\alpha\log(1-\lambda)+\beta h C_\lambda$, every point moves in the same
direction by at most
\[
 \frac\alpha8\left(\frac{\lambda^n}{1-\lambda}\right)^2
 =\frac\alpha8(e^{\beta h/\alpha}-1)^2
 \le\frac{\beta^2(503/500)^2}
 {160(1-\beta(503/500)/20)^2}.
\]
The first estimate follows from the second derivative of $\log(1+zs)$
and $s(1-s)\le1/4$; the last uses $e^x-1\le x/(1-x)$ for $0\le x<1$.
The bound decreases when $\alpha$ increases from twenty.
Moving points in one direction by at most this amount enlarges a circle
gap by at most the same amount: remove that amount from the appropriate
end of a purported larger empty arc and pull it back.

Compare this error with $91/200$ minus $503/500$ times the preceding
three-piece gap bound at $\lambda=183/500$.
On the first piece their ratio increases with $\beta$.
On the middle piece its derivative has the sign of
\[
 2\frac{91}{200}-4\frac{503}{500}
 +\frac{503}{500}\beta+\frac{(503/500)^2}{20}\beta^2,
\]
which increases but is still
$-381191095599/312500000000<0$ at its right endpoint $\beta=433/250$.
On the last piece the ratio increases.
At the two possible maxima, the available width minus the error is
\[
 \begin{array}{c|c}
 \beta=500/317&120809903739/10800370373000>0\\
 \beta=11/4&8981604362993/74248380562500>0.
 \end{array}
\]
Hence every circle gap of $\alpha A$ is strictly shorter than $b$ for
$\alpha\ge20$. Together with the preceding ranges, this proves
\eqref{eq:coarse-strict-link} for all real $\alpha>1$.

Finally fix $\alpha=p/q>1$ in lowest terms and the integer $j$ supplied by
\eqref{eq:coarse-strict-link}. For sufficiently large integers $N$, define
on $[1-\lambda,1]$
\[
 H_N(v):=\log\left(
 \frac{1-(1-\lambda^{j+pN}v)^\alpha}
 {\alpha\lambda^{j+pN}}\right).
\]
Then $\displaystyle\lim_{N\to\infty}H_N(v)=\log v$ in $C^1$.
The thickness of $H_N((1-\lambda)+\lambda C_\lambda)$ remains greater than
one: relative to $A$, the derivative ratio of this change of coordinates
has limit one. Its hull tends to $[-b,0]$.
By \eqref{eq:coarse-strict-link}, a fixed point of $\alpha A$ remains
strictly inside its translate by $\log\alpha-jh$.
This prevents that translated Cantor set from lying in a gap of $\alpha A$.
Conversely, the diameter $\alpha b$ of $\alpha A$ exceeds the diameter of
the translated set for all large $N$, so $\alpha A$ cannot lie in one of
its bounded gaps. The strict hull linkage excludes unbounded gaps as well.
The Gap Lemma yields $u,v\in(1-\lambda)+\lambda C_\lambda$ with
\[
 \alpha\log u-H_N(v)=\log\alpha-jh.
\]
Equivalently,
\[
 x:=\lambda^{qN}u,\qquad y:=1-\lambda^{j+pN}v
 \quad\hbox{satisfy}\quad x^\alpha+y^\alpha=1.
\]
For $j+pN\ge0$ these points belong to $C_\lambda^2$ and are nontrivial.
Their $x$-coordinates lie in the pairwise disjoint intervals
$\lambda^{qN}[1-\lambda,1]$, since $\lambda^q<1-\lambda$.
Thus the intersection contains infinitely many distinct nontrivial points.

\par\medskip\noindent\emph{Isolated points and their addresses.}
We first classify isolated points of a separable level curve with positive derivatives. A restriction to two Cantor cylinders, followed by affine rescaling, allows us to normalise its tangent ratio to $1\le r<1/\lambda$. Interchanging the coordinates and changing one cylinder depth replaces $r$ by $1/(\lambda r)$.
For the affine function $x+ry$, the four child hulls, in the order $00,10,01,11$, are
\[
[0,\lambda(1+r)],\quad
[1-\lambda,1+r\lambda],\quad
[r(1-\lambda),\lambda+r],\quad
[(1-\lambda)(1+r),1+r].
\]
If $1/\lambda-2<r<1/(1-2\lambda)$, consecutive hulls overlap strictly. The same holds throughout every sufficiently small descendant of a $C^1$ curve with this tangent ratio. Iterating the covering of each parent hull by its child hulls shows, by compactness, that every level in a descendant hull has a Cantor preimage there. Hence, along the address of an isolated point, only its actual child hull can contain the level once the depth is sufficiently large. Another child would give a distinct point arbitrarily close to it. This also excludes contacts with another child at a hull endpoint.

Suppose first that
\[
1\le r<\frac{\lambda}{1-2\lambda}.
\]
The exclusive regions of the mixed children $10$ and $01$ are empty, so an isolated address eventually uses only $00$ and $11$. At $r=1$ the two mixed hulls coincide. Under a small $C^1$ perturbation, any possible exclusive strips near their two ends are still covered by $00$ and $11$, respectively, since both affine covering margins are $3\lambda-1>0$. Write these paired symbols as $0$ and $3$. Their exclusive levels are at most $1-\lambda$ and at least $\lambda+r$, respectively. If $033$ occurred, the level after the first symbol would be at most $(1-\lambda)/\lambda$, whereas the next two symbols would require it to be at least
\[
(1-\lambda)(1+r)+\lambda(\lambda+r)
=1-\lambda+\lambda^2+r.
\]
These bounds are incompatible whenever
$r>1/\lambda-2+\lambda-\lambda^2$.
The expression on the right decreases on $[c,1/2)$ and equals $4c-1/2<1$ at $c$. Reflection excludes $300$. After the first change of symbol, the address must therefore alternate. The comparisons have positive margins, so the same exclusions hold after sufficiently small $C^1$ perturbations. Thus each coordinate address is eventually constant or alternating, including when $r=1$.

The child comparisons in the proof of Lemma~\ref{lem:two-bands}, before the constant mixed addresses are excluded, also apply throughout
\[
\max\{1,1/\lambda-1-\lambda\}<r<\frac1{1-\lambda-\lambda^2}.
\]
They force a continuing $10$ to be followed by $10$, and similarly for $01$, while $00$ cannot be followed by $11$ and conversely. An address avoiding overlap is consequently eventually constant. Again the strict comparisons persist under the vanishing linearisation errors. In either interval, an isolated point therefore has coordinates in $\mathbb Q(\lambda)$.

For $\lambda>c$ these intervals cover all normalised phases after the permitted coordinate exchange. Indeed
\[
\frac{\lambda}{1-2\lambda}>\frac1\lambda-1-\lambda,
\]
because the numerator of their difference is $3\lambda-1-2\lambda^3$, which vanishes at $c$ and increases on $[c,1/2]$. If $r\ge1/(1-\lambda-\lambda^2)$, then $1/(\lambda r)\le1/\lambda-1-\lambda$ and belongs to the first interval. Thus, for $\lambda>c$, every isolated positive point of the Fermat curve has eventually constant or alternating addresses and coordinates in $\mathbb Q(\lambda)$.

\par\medskip\noindent\emph{The two boundary phases at $\lambda=c$.}
At $c$, the intervals above cover every phase except $1+c$ and $2$. The second is reduced to the first by $r\mapsto1/(cr)$. We prove that an isolated point of a $C^2$ curve with algebraic coordinates and tangent ratio $1+c$ has coordinates in $\mathbb Q(c)$.
For this ratio, the limiting closed regions allowed by the unique-child condition are
\[
[0,1-c],\qquad \{1/2+c\},\qquad
\{3/2\},\qquad [1+2c,2+c].
\]
The two singletons are limiting exclusive regions; at the exact affine ratio their levels are covered by more than one child.
Let $x_n,y_n\in[0,1]$ denote its coordinate tails after $n$ paired digits, and define the tangent level
$z_n:=x_n+(1+c)y_n$.
Taylor expansion of the $C^2$ curve at the fixed point shows that $z_n$ is within $Kc^n$ of the exclusive set corresponding to its next paired digit, for some fixed $K>0$. Indeed every child-hull endpoint has this error after division by the first derivative and the cylinder scale; the unique-child condition proved above then gives the asserted bound.
If the next digits are $\varepsilon_{n+1},\eta_{n+1}\in\{0,1\}$, the tails satisfy the exact recurrence
\[
z_{n+1}
=\frac{z_n-(1-c)\bigl(\varepsilon_{n+1}+(1+c)\eta_{n+1}\bigr)}{c}.
\]
A mixed $10$ sends the centre $1/2+c$ to $1-c$, and a mixed $01$ sends $3/2$ to $1+2c$. The paired digits $00$ and $11$ exchange the latter two centres. Each centre is separated from every wrong exclusive set by at least $2c-1/2>0$.
Choose an integer $L\ge1$ with $2Kc^L<2c-1/2$. After a mixed digit at depth $n$, the next $n-L$ paired digits must alternate $00,11$: after $k$ shifts the initial error is at most $Kc^{n-k}$, whereas the exclusive-set error is at most $Kc^{n+k}$, and their sum is less than $2c-1/2$ for $k\le n-L$. Changing the depth convention only changes $L$ by one.

If mixed digits occur only finitely often, the strict $033$ and $300$ exclusions already prove eventual constancy or alternation. Otherwise each coordinate address has arbitrarily long alternating blocks whose lengths are at least their starting positions minus a fixed constant. Such a nonperiodic address satisfies the stammering condition of \cite[Theorem~5]{AdamczewskiBugeaud2007}: take an even word of approximately half the block length and repeat it twice. The preceding prefix length divided by this word length is bounded, and a subsequence makes the word lengths strictly increasing. A Cantor coordinate equals $(1/c-1)\sum_{j\ge1}\varepsilon_jc^j$; its digit-series value is algebraic whenever the coordinate is algebraic. Since $1/c=1+\sqrt3$ is a Pisot number, that theorem puts the digit-series value in $\mathbb Q(\sqrt3)$ unless it is transcendental. The algebraic-coordinate hypothesis excludes the latter alternative. Already periodic addresses also give values in this field.

For rational $\alpha>1$, a critical tangent equation makes $y/x$ algebraic, since $(y/x)^{\alpha-1}$ is an algebraic multiple of an integral power of $c$. The equation $x^\alpha+y^\alpha=1$ then makes both positive coordinates algebraic. Thus every positive isolated rational Fermat point at $c$ has coordinates in $\mathbb Q(c)$.

\par\medskip\noindent\emph{From a perfect fibre to strict returns.}
We give the dimension argument, including its stability when the contraction
changes. Fix $\lambda\in[c,1/2)$ and a positive ratio
$r\notin\mathbb Q(\lambda)$ in one of the open periodicity intervals
established above, after a fixed cylinder restriction and, if necessary,
interchanging coordinates. Put $L(s,t):=s+rt$.
The four child hulls of $L$ cover $[0,1+r]$ with strict consecutive
overlaps. Iteration and compactness therefore give
$L(C_\lambda^2)=[0,1+r]$.
The address argument above shows that every isolated point of every affine
fibre has coordinates in $\mathbb Q(\lambda)$.

A compact countable fibre containing two points would have two isolated
points: restrict it to two disjoint product cylinders containing the
original points, and use the existence of an isolated point in each
nonempty compact countable set. Product cylinders are relatively open
in $C_\lambda^2$, so these points are isolated in the whole fibre.
Equality of their affine levels would put $r$ in $\mathbb Q(\lambda)$,
a contradiction. Thus a fibre containing at least two points is uncountable.

Let $B$ be the union of \emph{all} pairwise intersections of the four
closed child hulls. This is a compact set. Every level in $B$ has two
distinct Cantor preimages, because each child hull is filled and distinct
product children are disjoint. Its fibre is therefore uncountable.
In an uncountable fibre, some descendant has two uncountable children.
Otherwise following the unique uncountable child at each step expresses
the fibre as one limiting point together with countably many countable
offshoots. At a node with two uncountable children, the level lies in the
interior of each of their hulls: an extreme level of a child has a unique
preimage, since both coefficients of $L$ are positive. The normalized
level consequently belongs to $\operatorname{int}B$.
Repeating this argument separately inside the two uncountable children
produces two incomparable nonempty words whose normalized levels return
to $\operatorname{int}B$. This works for every initial level in $B$ and
does not require the absence of triple overlaps.

For a product word $w$, let $z_w(\lambda)$ be its lower corner and write
\[
 \Phi_{w,\lambda}(u)
 :=\lambda^{-|w|}\bigl(u-L(z_w(\lambda))\bigr).
\]
The two strict return inequalities persist on a neighborhood of their
initial level. Compactness of $B$ gives finitely many word pairs, a maximum
length $M$, an open bounded neighborhood $U\supset B$, and $\eta>0$ such
that every $u\in U$ has a selected pair returning to points of $B$ at
distance at least $\eta$ from $\mathbb R\setminus B$.
Consider a normalized equation $L+e-u=0$, where $e(0,0)=0$ and
$\operatorname{Lip}(e)<\eta\lambda^M/2$, using the sup norm.
Restriction to $w$ replaces the error by
\[
 \lambda^{-|w|}
 \bigl(e(z_w+\lambda^{|w|}(s,t))-e(z_w)\bigr).
\]
Its Lipschitz constant does not increase, and the level correction is
bounded by $\lambda^{-|w|}\operatorname{Lip}(e)<\eta/2$.
Both children still return to $B$, so this construction iterates.

Every infinite branch converges to a zero, because the original residual
is the cylinder scale times a bounded normalized residual.
Assigning mass $2^{-j}$ to each binary generation-$j$ cylinder gives a
positive-dimensional measure. Indeed their depths are at most $jM$,
and any two distinct such cylinders are separated in at least one
coordinate by at least $(1-2\lambda)\lambda^{jM-1}$.
At radii comparable to $\lambda^{jM}$ a ball therefore meets a bounded
number of these cylinders. The mass distribution principle gives
dimension at least $\log2/(M|\log\lambda|)>0$.

Now suppose a $C^1$ level curve contains a compact perfect subset, and
at one of its points the normalized tangent ratio is the chosen $r$.
In the depth-$n$ product cylinder containing that point, subtract the
lower-corner value and divide by the first derivative at the point times
$\lambda^n$. The equation becomes $L+e_n-u_n=0$, with $e_n(0,0)=0$ and
$\displaystyle\lim_{n\to\infty}\operatorname{Lip}(e_n)=0$.
There are arbitrarily deep branching cylinders along this address;
otherwise the point would be isolated in the perfect subset.
At a branching cylinder, two distinct child image hulls of $L+e_n$
contain $u_n$. For the affine hulls, this implies
$\displaystyle\lim_{n\to\infty}\operatorname{dist}(u_n,B)=0$ along these branching depths:
there are only finitely many child
pairs, disjoint pairs have a fixed positive separation, and for an
intersecting pair the endpoint errors tend to zero.
A sufficiently deep branching cylinder thus has $u_n\in U$ and the
required Lipschitz bound. The preceding binary construction proves
positive Hausdorff dimension of the actual fibre.

\par\medskip\noindent\emph{Stability in the contraction parameter.}
The construction just obtained is also stable when $\lambda$ changes.
Keep its ratio $r$, compact core $B$, and finitely many return pairs fixed
at the original contraction. Each $z_w(\lambda)$ is a finite polynomial
in $\lambda$, so each selected map $\Phi_{w,\lambda}$ is continuous in
$\lambda$ and its level argument. Shrinking $U$ if necessary, the finite
cover and its positive margins persist for all contractions in a
neighborhood of the original one. The allowable Lipschitz bound can be
chosen uniformly there, since $M$ is fixed and the contractions stay
bounded away from zero and one half. The same restriction calculation
then constructs a binary tree for every nearby contraction.
The finite initial entrance cylinder also varies continuously with the
contraction. Its normalized level remains in $U$, and its defining
function remains within the chosen Lipschitz bound, after making the
neighborhood smaller. Thus positivity holds for the same actual level,
not merely for a nearby unspecified level.
In particular, a strict-return construction at $\lambda=c$ persists
for all $\lambda$ in some interval $(c-\varepsilon,c+\varepsilon)$.
No arithmetic classification at the perturbed contractions is needed.


\paragraph{Completion for rational exponents.}
Let $\alpha=p/q>1$ be reduced, with $p\ge4$. The endpoint construction
above gives infinitely many distinct nontrivial points for every
$c\le\lambda<1/2$. The address classification, including the boundary
phase argument at $c$, puts every positive isolated point in
$\mathbb Q(\lambda)^2$.
Its positive $q$th coordinate roots also belong to $\mathbb Q(\lambda)$.
Indeed, in a normal extension containing $x^{p/q}$ and $y^{p/q}$,
an automorphism over $\mathbb Q(\lambda)$ multiplies these two positive
numbers by $q$th roots of unity. Their sum remains one; equality in
the triangle inequality forces both roots of unity to equal one.
Thus $x^{p/q},y^{p/q}\in\mathbb Q(\lambda)$, and B\'ezout's identity
for $p,q$ gives $x^{1/q},y^{1/q}\in\mathbb Q(\lambda)$.
Each isolated point consequently gives a point on the projective
Fermat curve of degree $p$ over $\mathbb Q(\lambda)$. This curve is
smooth and has genus $(p-1)(p-2)/2\ge3$. If $\lambda$ is algebraic,
Faltings's theorem \cite{Faltings1983} gives only finitely many such
points. If $\lambda$ is transcendental, identify $\mathbb Q(\lambda)$
with $\mathbb Q(T)$. A nonconstant point would give a nonconstant
rational map from $\mathbb P^1$ to this smooth projective curve.
It extends to a morphism, which is impossible by Riemann--Hurwitz.
All coordinates therefore belong to $\mathbb Q$, and Faltings's
theorem again gives finiteness. Hence the full compact intersection
has only finitely many isolated points, including its possible axis
points.

Removing these finitely many isolated points leaves a nonempty perfect
compact set: the complement is closed, and a point that became isolated
after this removal would already have been isolated before it.
The ratio $(y/x)^{\alpha-1}$ is strictly monotone along
the positive Fermat arc. The field $\mathbb Q(\lambda)$ is countable,
and the set of cylinder endpoints is countable, so we can choose
a non-axis point of this perfect set whose ratio is
outside that field and whose coordinates are not cylinder endpoints.
Changing cylinder depths multiplies the ratio by an integral power of
$\lambda$, and coordinate exchange inverts it; these operations
preserve exclusion from $\mathbb Q(\lambda)$. The normalised ratio
lies in one of the open periodicity intervals above. At $\lambda=c$,
it cannot be either boundary phase, since $1+c$ and $2$ belong to
$\mathbb Q(c)$. The strict-return argument therefore gives a
positive-dimensional subset of the actual intersection.

This proves positivity for every exponent and every
$c\le\lambda<1/2$. At $\lambda=1/2$, the Cantor set is $[0,1]$ and
the Fermat arc has dimension one. For each fixed exponent the source
at $\lambda=c$ is described by finitely many strict returns and a
finite entrance cylinder, or by the strict inequalities of
Lemma~\ref{lem:two-bands}. The persistence proved above supplies
$\varepsilon>0$ such that positivity also holds for
$c-\varepsilon<\lambda<c$. Together with the whole upper interval,
this gives $\lambda_{\max}(\alpha)<c$. The value of $\varepsilon$
may depend on $\alpha$.
\end{proof}
